\chapter{KẾT LUẬN VÀ HƯỚNG PHÁT TRIỂN}

\section{Kết luận}

Đồ án đã xây dựng UniTrack như một nền tảng Web hỗ trợ giảng viên giám sát quá trình thực hiện dự án sinh viên. Hệ thống tập trung vào luồng dự án, nhiệm vụ, nộp bài và đánh giá, đồng thời cung cấp các chức năng nền tảng như xác thực, quản trị tài khoản, workspace, dashboard, tài nguyên và tệp minh chứng.

Thông qua quá trình thực hiện, đồ án đã áp dụng các kỹ thuật phát triển Web hiện đại, thiết kế API REST, quản lý trạng thái frontend, kiểm soát truy cập backend, ràng buộc cơ sở dữ liệu và kiểm thử tự động. Kết quả đạt được là một sản phẩm có thể tiếp tục mở rộng theo các lát tính năng rõ ràng.

\section{Hướng phát triển}

Trong tương lai, UniTrack có thể tiếp tục phát triển theo các hướng sau:

\begin{itemize}
  \item Hoàn thiện quản trị toàn bộ dự án cho vai trò quản trị viên.
  \item Xây dựng giao diện nhật ký hoạt động để theo dõi lịch sử thay đổi quan trọng.
  \item Chuyển lưu trữ tệp minh chứng sang dịch vụ bền vững như Cloudflare R2 hoặc S3-compatible storage.
  \item Bổ sung chính sách dung lượng, loại tệp, quét mã độc và sao lưu dữ liệu minh chứng.
  \item Mở rộng kiểm thử frontend cho các form dự án, nhiệm vụ, tài nguyên, tệp và luồng thêm/xóa thành viên.
  \item Tiếp tục cải thiện giao diện dashboard, trạng thái bài nộp và trải nghiệm duyệt bài của giảng viên.
\end{itemize}

Những hướng phát triển này giúp hệ thống tiến gần hơn tới một sản phẩm có thể sử dụng trong môi trường học tập thực tế với yêu cầu cao hơn về vận hành, an toàn dữ liệu và trải nghiệm người dùng.
